\section{Containerized Deployment} Following the architectural decisions described in Chapter~\ref{chap:system_architecture}, all components were deployed as Docker containers. Containerization provides reproducibility, isolation of dependencies, and simplifies the deployment and maintenance of the platform. 

\subsection{Docker Compose Organization} 
To improve maintainability and modularity, the platform is organized as a collection of independent Docker Compose files. Each architectural layer maintains its own deployment definition, while a main \texttt{docker-compose.yml} file aggregates all services using the Docker Compose \texttt{include} directive. 

This approach allows each subsystem to be developed and deployed independently while preserving a unified deployment strategy for the complete platform. 

Listing~\ref{lst:main-compose} shows the main Docker Compose configuration used to orchestrate the different services.

\begin{lstlisting}[language=yaml, caption={Main Docker Compose structure}, label={lst:main-compose}, basicstyle=\ttfamily\small, breaklines=true]
include:
  - docker/mqtt-broker/docker-compose.yml
  - docker/source_systems/vda5050_vehicle_simulator/docker-compose.yml
  ...
\end{lstlisting}

\subsection{Network Segmentation} 
Dedicated Docker bridge networks were created to isolate communication between system components. Rather than exposing every service to a common network, individual networks were assigned according to the communication requirements of each layer. 

This strategy reduces unnecessary connectivity between containers and simplifies service discovery. Furthermore, it follows the principle of least privilege by restricting access only to the services required by each component. 

The current implementation defines dedicated networks for MQTT communication, event streaming, and data processing services.

Listing~\ref{lst:docker-networks} presents the network configuration currently used by the platform.

\begin{lstlisting}[language=yaml, caption={Docker network configuration}, label={lst:docker-networks}, basicstyle=\ttfamily\small, breaklines=true]
networks:
  spark-net:
    driver: bridge
  redpanda-net:
    driver: bridge
  mqtt-net:
    driver: bridge
\end{lstlisting}
   
